\section{Summary of the Work}

This project developed \systemname, a lightweight multi-modal framework for privacy-constrained academic assistance. The system separates public and protected corpora, ingests PDF/image/text with metadata, classifies Bloom levels with a fine-tuned Qwen LoRA model, answers student questions with Qwen GGUF RAG, and screens outputs with PrivacyGuard.

\section{Key Findings and Contributions}

\begin{itemize}
    \item \textbf{Qwen LoRA Bloom:} 0.748 accuracy, 0.721 macro-F1, 0.880 within-one-level accuracy on the official test split ($n{=}2330$).
    \item \textbf{Baselines:} TF-IDF + LinearSVC 0.839 accuracy on hold-out; zero-shot Qwen GGUF 0.441---LoRA fine-tuning is required for reliable moderation labels.
    \item \textbf{Privacy:} 100\% block rate on 104 student reconstruction attacks; 86.7\% benign allow rate.
    \item \textbf{Qwen RAG:} 0.914 token-F1, 0.800 exact match, 1.000 retrieval hit@3 on the local QA benchmark.
    \item \textbf{OCR/multimodal:} 0.963 OCR token-F1 on synthetic typed images; smoke tests passed.
\end{itemize}

\section{Limitations of the Study}

Synthetic privacy prompts, small QA benchmark, hold-out vs.\ official-test split difference for Bloom tables, and simulated federation. No formal differential privacy.

\section{Recommendations and Future Work}

Expand red-teaming, larger course QA sets, full hold-out LoRA evaluation on the same 379-item split, handwriting OCR, and production federated deployment.

Overall, \systemname{} shows that Qwen LoRA Bloom labeling plus role-separated RAG is a practical university prototype.
